% TEMPLATE-MILC-v3.tex - plantilla maestra para todas las guias MILC v3
% Ver scripts/generadores/build-guia-milc.py para la lista completa de
% claves esperadas. El script valida que no queden placeholders sin
% reemplazar antes de escribir el archivo final.

\documentclass[11pt,letterpaper]{article}
\usepackage[margin=1.75cm]{geometry}
\usepackage[table]{xcolor}
\usepackage{fontspec}
\usepackage[spanish]{babel}
\usepackage{tikz}
\usepackage[most]{tcolorbox}
\usepackage{tabularx}
\usepackage{array}
\usepackage{fancyhdr}
\usepackage{titlesec}
\usepackage{ragged2e}
\usepackage{enumitem}
\usepackage{microtype}

\setmainfont{Helvetica Neue}
\setsansfont{Helvetica Neue}
\setlength{\parindent}{0pt}
\setlength{\parskip}{6pt}
\hyphenpenalty=200
\exhyphenpenalty=200
\pretolerance=400
\tolerance=2000
\setlength{\emergencystretch}{1em}
\renewcommand{\arraystretch}{1.25}
\setlist{leftmargin=5mm,itemsep=1mm,topsep=1mm}
\newcolumntype{Y}{>{\RaggedRight\arraybackslash}X}

\definecolor{milcMagenta}{HTML}{E6007E}
\definecolor{milcVino}{HTML}{5A0038}
\definecolor{milcTurquesa}{HTML}{0093A5}
\definecolor{milcVerde}{HTML}{58A923}
\definecolor{milcMostaza}{HTML}{E5B400}
\definecolor{milcOcre}{HTML}{B86600}
\definecolor{milcNegro}{HTML}{241816}
\definecolor{milcGris}{HTML}{F7F7F4}
\definecolor{milcLinea}{HTML}{E8E2DD}
\definecolor{milcGradePrimary}{HTML}{7C3AED}
\definecolor{milcGradeDark}{HTML}{4C1D95}
\definecolor{milcGradeSoft}{HTML}{EDE9FE}
\definecolor{milcGradeAccent}{HTML}{CFE7FF}
\definecolor{milcGradeText}{HTML}{FFFFFF}
\color{milcNegro}

\pagestyle{fancy}
\fancyhf{}
\lhead{\small Tecnología e Informática · Grado 9}
\rhead{\small Guía 13 · MILC v3}
\cfoot{\small\thepage}
\renewcommand{\headrulewidth}{0pt}

\newtcolorbox{titlebox}[1]{
  enhanced,colback=#1,colframe=#1,arc=7mm,boxrule=0pt,
  left=7mm,right=7mm,top=3mm,bottom=3mm,
  before skip=8pt,after skip=8pt,
  fontupper=\sffamily\bfseries\Large\color{white}
}

\newtcolorbox{softbox}[3]{
  enhanced,colback=#2,colframe=#1,borderline west={4pt}{0pt}{#1},
  arc=2mm,boxrule=.4pt,left=4mm,right=4mm,top=3mm,bottom=3mm,
  before skip=6pt,after skip=8pt,title={#3},coltitle=white,
  colbacktitle=#1,fonttitle=\sffamily\bfseries,fontupper=\RaggedRight,
  attach boxed title to top left={xshift=3mm,yshift=-2mm},
  boxed title style={arc=2mm, boxrule=0pt}
}

\newtcolorbox{drawbox}[2]{
  enhanced,colback=white,colframe=#1,arc=4mm,boxrule=1pt,
  left=5mm,right=5mm,top=4mm,bottom=4mm,before skip=8pt,after skip=8pt,
  title={#2},coltitle=white,colbacktitle=#1,fonttitle=\sffamily\bfseries,
  fontupper=\RaggedRight,
  attach boxed title to top left={xshift=4mm,yshift=-2mm},
  boxed title style={arc=3mm, boxrule=0pt}
}

\newtcolorbox{infoband}[1]{
  enhanced,colback=#1!8,colframe=#1!50,arc=2mm,boxrule=.5pt,
  left=4mm,right=4mm,top=2mm,bottom=2mm,before skip=4pt,after skip=4pt,
  fontupper=\small\RaggedRight
}

% Puente narrativo entre secciones (contrato editorial Conéctate).
% Da continuidad: "Acabas de X. Ahora vas a Y porque Z."
% Visualmente: banda izquierda en color del grado, texto en itálica pequeña.
\newtcolorbox{puentebox}{
  enhanced,colback=milcGradeSoft,colframe=milcGradeDark,
  borderline west={3pt}{0pt}{milcGradeDark},
  arc=0mm,boxrule=0pt,left=5mm,right=4mm,top=2.5mm,bottom=2.5mm,
  before skip=4pt,after skip=8pt,
  fontupper=\itshape\small\color{milcNegro}\RaggedRight
}
\newcommand{\puente}[1]{%
  \begin{puentebox}%
    \textcolor{milcGradeDark}{\textbf{\textrightarrow}}\hspace{2mm}#1%
  \end{puentebox}%
}

\newcommand{\linead}[1]{\noindent\color{milcLinea}\rule{#1}{0.5pt}\color{milcNegro}}
\newcommand{\respuesta}[1]{\par\vspace{2mm}\foreach \n in {1,...,#1}{\noindent\color{milcLinea}\rule{\linewidth}{0.5pt}\par\vspace{5mm}}\color{milcNegro}}
\newcommand{\checkbox}{\textcolor{milcGradeDark}{\fbox{\phantom{x}}}\hspace{5pt}}

\newcommand{\phasecard}[4]{
\begin{tcolorbox}[enhanced,colback=#3,colframe=#2,arc=3mm,boxrule=.5pt,height=3.05cm,valign=center,halign=center,left=1.5mm,right=1.5mm,top=2mm,bottom=2mm]
{\sffamily\bfseries\fontsize{22}{24}\selectfont\textcolor{#2}{#1}}\par
{\sffamily\bfseries\small #4}
\end{tcolorbox}
}

\begin{document}
\thispagestyle{empty}
\begin{tikzpicture}[remember picture,overlay]
  \fill[white] (current page.south west) rectangle (current page.north east);
  \path[fill=milcGradePrimary,rounded corners=16mm]
    ([xshift=.55cm,yshift=-.55cm]current page.north west) rectangle
    ([xshift=-.55cm,yshift=3.65cm]current page.south east);

  \node[anchor=north west,text=milcGradeText] at ([xshift=2.0cm,yshift=-1.85cm]current page.north west)
    {\sffamily\bfseries\fontsize{14}{16}\selectfont GRADO};
  \node[anchor=north west,text=milcGradeText,opacity=.82] at ([xshift=2.0cm,yshift=-2.75cm]current page.north west)
    {\sffamily\bfseries\fontsize{9}{11}\selectfont SERIE GUÍAS · TECNOLOGÍA E INFORMÁTICA};

  \begin{scope}[shift={([xshift=-3.05cm,yshift=-2.55cm]current page.north east)}]
    \fill[white,opacity=.80] (-.62,-.52) rectangle (.62,.52);
    \draw[milcGradeText,line width=1pt] (-.62,-.52) rectangle (.62,.52);
    \fill[milcGradeDark] (-.42,-.38) rectangle (-.22,.06);
    \fill[milcGradeAccent] (-.10,-.38) rectangle (.10,.24);
    \fill[milcGradePrimary!60!black] (.22,-.38) rectangle (.42,.38);
  \end{scope}

  \node[anchor=west,text=milcGradeText] at ([xshift=1.9cm,yshift=-12.85cm]current page.north west)
    {\sffamily\bfseries\fontsize{124}{124}\selectfont 9};
  \draw[milcGradeText,line width=1.7pt]
    ([xshift=4.52cm,yshift=-11.68cm]current page.north west) circle (.13cm);

  \node[anchor=west,text=milcGradeText,text width=8.7cm,align=left] at ([xshift=10.35cm,yshift=-11.6cm]current page.north west)
    {
     {\sffamily\bfseries\fontsize{19}{22}\selectfont Tipografía ---\\la voz visible\\de las palabras}\\[4mm]
     {\sffamily\bfseries\fontsize{23}{26}\selectfont Guía 13-9-TIC}\\[2mm]
     {\sffamily\fontsize{13}{16}\selectfont Período 2 · Diseño editorial digital}\\[5mm]
     {\sffamily\bfseries\fontsize{11}{13}\selectfont Institución Educativa Sor María Juliana}\\[1mm]
     {\sffamily\fontsize{10}{12}\selectfont Autor: Dr. Álvaro Cárdenas Orozco}
    };

  \path[fill=milcGradeSoft,rounded corners=7mm]
    ([xshift=1.35cm,yshift=.85cm]current page.south west) rectangle
    ([xshift=-1.35cm,yshift=4.25cm]current page.south east);
  \draw[milcGradeDark,line width=.65pt,rounded corners=7mm]
    ([xshift=1.35cm,yshift=.85cm]current page.south west) rectangle
    ([xshift=-1.35cm,yshift=4.25cm]current page.south east);
  \node[anchor=center,text=milcNegro,text width=17.0cm,align=center] at ([yshift=2.55cm]current page.south)
    {\sffamily\fontsize{10}{12}\selectfont Metodología MILC · Apertura ancestral · Pensamiento computacional · Triángulo Dussel-Estoicismo-Floridi\\
    \textcolor{milcGradeDark}{\bfseries Producto:} Clasificación de 10 tipografías (display vs cuerpo, con/sin serifa, geométrica/humanista) + par tipográfico propio elegido para la revista del periodo, justificado en 1 frase};
\end{tikzpicture}
\mbox{}
\newpage

\section*{Apertura: Tipografía --- la voz visible de las palabras}

\begin{softbox}{milcGradeDark}{milcGradeSoft}{Saber ancestral}
Antes de la tipografía digital existió la \textbf{caligrafía}: los \emph{rótulos de barberías y panaderías} del Valle pintados a mano con elegancia, las \emph{letras de las décimas} del Pacífico escritas con tinta sobre cuaderno, la \emph{escritura cuidada} de los registros parroquiales, las \emph{cartas de amor} de tu abuela con su letra inconfundible. Cada tipografía digital hereda gestos de mano humana: una letra con serifa serenamente curvada viene del trazo de la pluma sobre pergamino. Una letra geométrica viene de la regla y el compás del XX. Las tipografías cargan siglos de oficio.
\end{softbox}

\begin{softbox}{milcTurquesa}{milcGris}{Saber contemporáneo}
Hoy elegir tipografía es un acto editorial mayor: define la \emph{voz visible} de tus palabras. Una tipografía con serifa (Times, Garamond, Lora) se siente clásica, autorizada. Una sin serifa (Helvetica, Inter, Roboto) se siente moderna, neutra. Una display (Bricolage, Playfair) se siente dramática, llamativa. Tu próxima revista necesitará un \textbf{par tipográfico}: una para titulares (display) y una para cuerpo (legible). Esa elección moldea cómo el lector recibe tu mensaje antes de leerlo.
\end{softbox}

\begin{softbox}{milcMostaza}{milcGris}{Pregunta-puente}
\textit{¿Por qué la letra de tu abuela en una carta antigua se ``ve cálida''? ¿Y por qué algunas tipografías digitales transmiten lo mismo y otras no?}
\end{softbox}

\begin{softbox}{milcVerde}{milcGris}{Saber y saber hacer}
Al terminar podrás: (1) \textbf{identificar} 10 tipografías y clasificarlas en 4 categorías; (2) \textbf{analizar} qué emoción transmite cada categoría; (3) \textbf{aplicar} la regla del ``par tipográfico'' (display + cuerpo); (4) \textbf{crear} tu par tipográfico propio para la revista del periodo, justificado.
\end{softbox}

\small
\begin{tabularx}{\textwidth}{>{\bfseries}p{3.0cm}Y}
\rowcolor{milcGradePrimary!12}Autor & Dr. Álvaro Cárdenas Orozco\\
DBA & Producción de piezas editoriales digitales con criterio estético y ético (MEN, T\&I)\\
Referentes & Tipografía colombiana · Diseño centrado en accesibilidad · Floridi\\
Producto final & Clasificación de 10 tipografías (display vs cuerpo, con/sin serifa, geométrica/humanista) + par tipográfico propio elegido para la revista del periodo, justificado en 1 frase\\
\end{tabularx}
\normalsize

\puente{Antes de empezar, mira el viaje completo. La sesión 2 te dio la grilla; hoy le pones \emph{voz} a esa grilla con tipografía. Sin tipografía bien elegida, hasta la mejor grilla se siente fría.}

\begin{titlebox}{milcGradeDark}
Ruta MILC: así vamos a aprender
\end{titlebox}
\begin{tabularx}{\textwidth}{>{\centering\arraybackslash}X>{\centering\arraybackslash}X>{\centering\arraybackslash}X>{\centering\arraybackslash}X}
\phasecard{1}{milcVerde}{milcGris}{Escuta\\[-1mm]\small Observo, escucho y pregunto.} &
\phasecard{2}{milcTurquesa}{milcGris}{Sistematización\\[-1mm]\small Pensamiento computacional.} &
\phasecard{3}{milcMagenta}{milcGris}{Praxis\\[-1mm]\small Construyo y aplico.} &
\phasecard{4}{milcVino}{milcGris}{Evaluación\\[-1mm]\small Reflexión liberadora.}
\end{tabularx}


\puente{Antes de elegir tu par tipográfico, vamos a entrenar el ojo en 10 tipografías reales. Cada una tiene historia, propósito y emoción. Aprender a leerlas es leer el ADN del diseño.}

\begin{titlebox}{milcVerde}
Fase 1 · Escuta
\end{titlebox}

\begin{softbox}{milcVerde}{milcGris}{Escena para escuchar}
\textbf{Actividad 1 · IDENTIFICA --- 10 tipografías reales} (15 min · individual). Busca \textbf{10 tipografías} en lugares cercanos a ti: 3 de libros (mira la portada y el cuerpo), 3 de revistas o periódicos, 2 de logotipos de marcas que conoces, 2 de letreros o rótulos físicos cercanos. Para cada una: ¿con serifa o sin serifa? ¿display (grande, llamativa) o cuerpo (legible para texto)? ¿qué emoción transmite?
\end{softbox}

\checkbox Encontré 10 tipografías reales (3 libros + 3 revistas + 2 logos + 2 rótulos).\\
\checkbox Identifiqué para cada una si es display o cuerpo, con o sin serifa.\\
\checkbox Anoté la emoción que transmite (clásica, moderna, dramática, neutra, etc.).

\begin{infoband}{milcVerde}
\textbf{Lo que va al cuaderno (Actividad 1 · IDENTIFICA).} Título: «Actividad 1 --- 10 tipografías bajo la lupa». \textbf{Formato:} tabla de 5 columnas (Tipografía o muestra~|~Origen~|~Display/Cuerpo~|~Con/Sin serifa~|~Emoción transmitida), 10 filas. \textbf{Sabes que terminaste cuando} reconoces patrones (por ejemplo, las display suelen ser llamativas, las de cuerpo legibles).
\end{infoband}

\puente{Ya viste 10 tipografías. Ahora vamos a entender las \textbf{4 categorías} en que se organizan: display vs cuerpo, con vs sin serifa, geométrica vs humanista. Combinarlas con criterio es el oficio del diseñador.}

\begin{titlebox}{milcTurquesa}
Fase 2 · Sistematización con pensamiento computacional
\end{titlebox}

\begin{softbox}{milcTurquesa}{milcGris}{Cuatro pilares aplicados al tema}
Toda tipografía se clasifica en 4 dimensiones cruzadas. Vamos a descomponerlas con los pilares del pensamiento computacional para que sepas combinar con criterio en lugar de gusto.
\end{softbox}

\small
\begin{tabularx}{\textwidth}{>{\bfseries\RaggedRight\arraybackslash}p{3.6cm}Y}
\rowcolor{milcTurquesa!90}\color{white}Pilar & \color{white}\bfseries Aplicación al tema\\
1. Descomposición & La clasificación tipográfica se descompone en 4 ejes: (1) \textbf{display vs cuerpo} (display para titulares grandes y dramáticos; cuerpo para texto largo legible), (2) \textbf{con vs sin serifa} (con serifa: pies en las letras, sensación clásica y de lectura larga; sin serifa: limpia, moderna, ideal pantalla), (3) \textbf{geométrica vs humanista} (geométrica: trazos uniformes basados en círculos perfectos; humanista: trazos con variación que evocan la mano), (4) \textbf{regular vs variable} (regular: pesos fijos; variable: un solo archivo con muchos pesos posibles).\\
2. Reconocimiento de patrones & Toda combinación fuerte sigue el patrón \textbf{``par contrastado, no idéntico''}: tu titular y tu cuerpo deben sentirse parientes pero distintos. Mismo concepto (humanista + humanista) puede aburrir. Concepto opuesto (display geométrica + cuerpo humanista con serifa) crea contraste rico.\\
3. Abstracción & Lo esencial cabe en una pregunta: \emph{¿qué voz quiero darle a mi revista y qué par tipográfico transmite esa voz?}. Sin esa pregunta, eliges Inter por defecto y la revista se ve genérica.\\
4. Algoritmo conceptual & Pasos para elegir tu par: (1) define emoción central de tu revista; (2) elige display que transmita esa emoción para titulares; (3) elige cuerpo que respete la emoción pero priorice legibilidad; (4) verifica contraste entre ambas; (5) prueba con un titular y un párrafo de cuerpo reales.\\
\end{tabularx}
\normalsize

\begin{drawbox}{milcTurquesa}{Anatomía de un par tipográfico}
\textbf{Display (titulares):} llamativa, con personalidad, 36pt+. Ej: Playfair, Bricolage, Oswald. \\[1mm] \textbf{Cuerpo (texto largo):} legible, neutra, 10-12pt. Ej: Inter, Source Sans, Lora. \\[1mm] \textbf{Contraste:} las dos deben sentirse parientes pero distintas. \\[1mm] \textbf{Pesos:} cada una en al menos 2 pesos (regular + bold). \\[1mm] \textbf{Coherencia con la marca:} las dos deben sostener la emoción central de la revista. \\[1mm] \textbf{Accesibilidad:} legibles a tamaños pequeños, en pantalla y papel.
\end{drawbox}

\begin{softbox}{milcMostaza}{milcGris}{Errores comunes y su corrección}
(1) Usar la misma tipografía para titulares y cuerpo (aburrido y sin jerarquía). (2) Usar dos display al mismo tiempo (caos visual). (3) Elegir tipografía decorativa para cuerpo (ilegible en textos largos). (4) Importar tipografías de moda sin probar legibilidad. (5) Olvidar que hay tipografías Google Fonts gratuitas excelentes --- no necesitas pagar.
\end{softbox}

\begin{infoband}{milcTurquesa}
\textbf{Lo que va al cuaderno (Actividad 2 · ANALIZA).} Título: «Actividad 2 --- 4 ejes de clasificación tipográfica». \textbf{Formato:} 4 fichas, una por eje, cada una con: definición + 2 tipografías que ejemplifican cada extremo del eje. \textbf{Sabes que terminaste cuando} puedes mirar una tipografía y ubicarla en los 4 ejes sin dudar.
\end{infoband}

\puente{Tienes el método. Toca decidir: tu par tipográfico para la revista del periodo. Una para titulares, una para cuerpo, ambas eligidas con justificación.}

\begin{titlebox}{milcMagenta}
Fase 3 · Praxis con pensamiento computacional
\end{titlebox}

\begin{softbox}{milcMagenta}{milcGris}{Construyendo el producto con los cuatro pilares}
\textbf{Actividad 3 · CREA --- Tu par tipográfico para la revista} (30 min · individual). Hora de elegir. Vas a definir el par tipográfico (display + cuerpo) que usarás durante todo el periodo 2 para tu revista digital. Esta elección será tu compromiso editorial.
\end{softbox}

\small
\begin{tabularx}{\textwidth}{>{\bfseries\RaggedRight\arraybackslash}p{3.6cm}Y}
\rowcolor{milcMagenta!90}\color{white}Pilar & \color{white}\bfseries Aplicación al producto\\
Descomposición & Tu par se descompone en 4 decisiones: (1) emoción central de tu revista (calma, urgencia, ternura, rigor, juego), (2) tipografía display elegida con su nombre exacto, (3) tipografía cuerpo elegida con su nombre exacto, (4) justificación de 1 frase por cada elección.\\
Patrones & Sigue el patrón \textbf{``Google Fonts antes que ilegales''}: usa Google Fonts (gratuitas y profesionales). Bricolage Grotesque, Inter, Lora, Playfair Display, Source Sans, IBM Plex Sans --- todas profesionales, todas gratis.\\
Abstracción & Cada elección expresa una sola idea. No pongas ``elegí Playfair porque es bonita y elegante y moderna y profesional'' --- elige UN adjetivo y defiéndelo.\\
Algoritmo de armado & Algoritmo: (1) declara emoción central; (2) ve a fonts.google.com y prueba 3 display que la transmitan; (3) ve a fonts.google.com y prueba 3 cuerpo que respeten la emoción pero priorizen legibilidad; (4) elige UNA de cada; (5) escribe en cuaderno un titular y un párrafo cortos usándolas (puedes pegar pantallazos o escribirlas a mano imitándolas).\\
\end{tabularx}
\normalsize

\begin{drawbox}{milcMagenta}{Checklist antes de cerrar el par tipográfico}
\checkbox Emoción central de la revista declarada en 1 palabra. \\[1mm] \checkbox Tipografía display elegida con nombre exacto. \\[1mm] \checkbox Tipografía cuerpo elegida con nombre exacto. \\[1mm] \checkbox Justificación de 1 frase para cada elección. \\[1mm] \checkbox Ambas son de Google Fonts (verificado en fonts.google.com). \\[1mm] \checkbox Probé un titular y un párrafo reales con ambas.
\end{drawbox}

\begin{softbox}{milcMagenta}{milcGris}{Plantilla de trabajo}
\textbf{Emoción central de mi revista.} \textit{1 palabra.} \\[1mm] \textbf{Display elegida.} \textit{Nombre + justificación de 1 frase.} \\[1mm] \textbf{Cuerpo elegida.} \textit{Nombre + justificación de 1 frase.} \\[1mm] \textbf{Prueba real.} \textit{Un titular y un párrafo usando ambas.}
\end{softbox}

\begin{infoband}{milcMagenta}
\textbf{Lo que va al cuaderno (Actividad 3 · CREA).} Título: «Actividad 3 --- Mi par tipográfico para la revista». \textbf{Formato:} emoción central + par tipográfico justificado + prueba real (titular y párrafo). \textbf{Extensión:} 1 página de cuaderno. \textbf{Sabes que terminaste cuando} puedes defender por qué ese par y no otro frente a un compañero.
\end{infoband}

\puente{Lo que sigue resume qué entregas hoy: clasificación de 10 + par tipográfico propio. El criterio: que tu par sirva al propósito de tu revista, no que ``se vea bonito''.}

\begin{titlebox}{milcMostaza}
Producto de la sesión
\end{titlebox}
\textbf{Par tipográfico (display + cuerpo) elegido y justificado para la revista del periodo}.
\begin{softbox}{milcMostaza}{milcGris}{Criterios de calidad}
(1) Emoción central de la revista declarada. (2) Display elegida con nombre exacto + justificación. (3) Cuerpo elegida con nombre exacto + justificación. (4) Ambas son de Google Fonts. (5) Hay prueba real con un titular y un párrafo. (6) El par mantiene contraste sin perder coherencia.
\end{softbox}

\puente{Antes de cerrar, mira la tipografía desde las cinco dimensiones humanas. Elegir tipografía es elegir voz; elegir voz es elegir a quién hablas y cómo lo respetas.}

\begin{titlebox}{milcVino}
Fase 4 · Evaluación liberadora — cinco dimensiones
\end{titlebox}

\begin{softbox}{milcVino}{milcGris}{Sentido de la evaluación}
Evaluar no es solo revisar una nota. Es reconocer qué aprendí, qué cambió en mí, qué emociones manejé, cómo me relaciono con la ciudad y la comunidad, y qué le dejaré a quien viene detrás.
\end{softbox}

\textbf{Mi reflexión liberadora en cinco dimensiones.} Completa cada frase iniciada:

\small
\begin{tabularx}{\textwidth}{>{\bfseries\RaggedRight\arraybackslash}p{3.4cm}Y>{\RaggedRight\arraybackslash}p{4.6cm}}
\rowcolor{milcVino}\color{white}Dimensión & \color{white}\bfseries Mi respuesta (frame starter) & \color{white}\bfseries Algo que descubrí\\
1. Personal & Lo que más cambió en mí fue \linead{6.5cm} & \linead{4.2cm}\\
2. Emocional & La emoción más fuerte fue \linead{2.5cm} y la regulé \linead{3cm} & \linead{4.2cm}\\
3. Ciudadana & En democracia esto importa porque \linead{6cm} & \linead{4.2cm}\\
4. Local (Cartago) & En mi barrio sería útil para \linead{6.5cm} & \linead{4.2cm}\\
5. Intergeneracional & A mi abuela/abuelo le diría \linead{6.5cm} & \linead{4.2cm}\\
\end{tabularx}
\normalsize

\begin{infoband}{milcVino}
\textbf{Para la próxima sustentación.} Vuelve a esta tabla en seis meses. Verás que las respuestas cambian — no porque lo aprendido cambie, sino porque tú cambias. La evaluación liberadora es una conversación contigo a través del tiempo.
\end{infoband}


\puente{Y para terminar, tres voces sobre la voz visible. Dussel sobre las tipografías colonizadas, Marco Aurelio sobre la simplicidad de las letras que duran, Floridi sobre la tipografía como infraestructura cognitiva.}

\begin{titlebox}{milcGradeDark}
Triángulo de pensamiento · cierre filosófico
\end{titlebox}

\begin{softbox}{milcVino}{milcGris}{Dussel · lente del nosotros}
\textit{``Toda tipografía habla con acento: las ``neutras'' son las que mejor disimulan su origen.''}

Las tipografías que se nos venden como ``neutras'' (Helvetica, Times) no lo son: vienen de la modernidad europea del siglo XIX y XX. Reconocer eso no es rechazarlas --- es saber qué eliges cuando las usas. Hoy existen tipografías latinoamericanas (Rambla, Andika, Source Sans con caracteres adicionales) y africanas (Source Sans Hebrew, NotoSans para múltiples scripts) que abren la voz visible de la infoesfera.

\textbf{¿Las tipografías que uso reflejan mi región y mi lengua o solo replican estándares globales? ¿Conozco tipografías hechas en o desde Latinoamérica?}
\end{softbox}
\linead{\linewidth}\\[1mm]
\linead{\linewidth}

\begin{softbox}{milcMostaza}{milcGris}{Estoicismo · Marco Aurelio · cuidado interior}
\textit{``Las letras que duran son las que respetan los gestos antiguos de la mano.''}

Las tipografías humanistas con serifa (Garamond, Caslon) llevan 400 años funcionando porque conservan los gestos del calígrafo medieval. Las fuentes geométricas de moda duran décadas. La sabiduría estoica del tipógrafo: elegir letras que han demostrado durar antes que letras que demuestren novedad.

\textbf{¿Elijo tipografías porque están de moda o porque tienen historia que demuestre que duran? ¿Qué cambia si priorizo la durabilidad?}
\end{softbox}
\linead{\linewidth}\\[1mm]
\linead{\linewidth}

\begin{softbox}{milcTurquesa}{milcGris}{Floridi · lente de la infoesfera}
\textit{``La tipografía es la infraestructura cognitiva más invisible y más poderosa de la infoesfera.''}

Cada artículo, cada email, cada notificación que lees está mediada por decisiones tipográficas que alguien tomó. Las decisiones de Helvetica del XX modelaron el siglo de la modernidad. Las decisiones de Inter del XXI están modelando el siglo digital. La tipografía no es decoración --- es infraestructura cognitiva colectiva.

\textbf{¿Qué tipografías han modelado mi forma de leer sin que lo note? ¿Si pudiera elegirlas yo, qué cambiaría?}
\end{softbox}
\linead{\linewidth}\\[1mm]
\linead{\linewidth}

\begin{titlebox}{milcVino}
Compromiso final
\end{titlebox}

\small
\begin{tabularx}{\textwidth}{>{\RaggedRight\arraybackslash}p{3.0cm}YYY}
\rowcolor{milcVino}\color{white}\bfseries Criterio & \color{white}\bfseries Lo logré & \color{white}\bfseries En proceso & \color{white}\bfseries Necesito apoyo\\
Comprensión del tema & Explico ideas centrales con mis palabras. & Reconozco ideas pero falta precisión. & Necesito repasar conceptos base.\\
Pensamiento computacional & Apliqué los 4 pilares al diseñar mi producto. & Apliqué algunos pilares. & Necesito releer la fase 2.\\
Aplicación y producto & Mi producto cumple los criterios y se puede revisar. & Producto funcional con ajustes pendientes. & Aún en construcción.\\
Reflexión 5 dimensiones & Respondí cada dimensión con honestidad. & Respondí algunas con detalle. & Mis respuestas son muy generales.\\
Triángulo de pensamiento & Conecté las 3 voces con mi trabajo real. & Conecté al menos una. & No relaciono las voces con mi proceso.\\
\end{tabularx}
\normalsize

\textbf{Hoy entendí que:}\\[1mm]
\linead{\linewidth}\\[1mm]
\linead{\linewidth}

\textbf{Mi compromiso para la próxima sesión es:}\\[1mm]
\linead{\linewidth}\\[1mm]
\linead{\linewidth}

\begin{softbox}{milcMostaza}{milcGris}{Cita de despedida · Marco Aurelio}
\textit{``El obstáculo en el camino se vuelve el camino. Lo que estorba al camino se vuelve camino.''}\\[1mm]
Cada error de hoy, bien observado, se convierte en pieza del camino que pisarás mañana.
\end{softbox}

\small
\begin{tabularx}{\textwidth}{>{\bfseries}p{3.6cm}Y}
\rowcolor{milcGradeDark!12}Nombre completo: & \linead{10cm}\\
Fecha: & \linead{4cm} \quad Curso/grupo: \linead{4cm}\\
Mi firma: & \linead{9cm}\\
\end{tabularx}
\normalsize

\begin{infoband}{milcGradeDark}
\textbf{Para la próxima semana.} Configura el par tipográfico que elegiste en tu próximo trabajo del colegio (puede ser en Google Docs, Word, Figma o Canva). Úsalo durante al menos 3 días. Anota qué notaste --- si te incomodó, si te gustó, si lo cambiarías. El próximo lunes traemos esa bitácora. La elección no es definitiva: puedes ajustarla.
\end{infoband}

\end{document}
